% =====================================================================
%  Week 1 Implementation — drop-in chapter section for the dissertation
%  Project: Adaptive Retrieval Gating for Safety-Aware Medical RAG
%  Author:  Tharit Mohamed Hussain Khan (22058691), KCL MSc AI, 7CCSMPRJ
%
%  This file is self-contained for the listings/booktabs packages it uses.
%  To include it in the main report:  % =====================================================================
%  Week 1 Implementation — drop-in chapter section for the dissertation
%  Project: Adaptive Retrieval Gating for Safety-Aware Medical RAG
%  Author:  Tharit Mohamed Hussain Khan (22058691), KCL MSc AI, 7CCSMPRJ
%
%  This file is self-contained for the listings/booktabs packages it uses.
%  To include it in the main report:  % =====================================================================
%  Week 1 Implementation — drop-in chapter section for the dissertation
%  Project: Adaptive Retrieval Gating for Safety-Aware Medical RAG
%  Author:  Tharit Mohamed Hussain Khan (22058691), KCL MSc AI, 7CCSMPRJ
%
%  This file is self-contained for the listings/booktabs packages it uses.
%  To include it in the main report:  % =====================================================================
%  Week 1 Implementation — drop-in chapter section for the dissertation
%  Project: Adaptive Retrieval Gating for Safety-Aware Medical RAG
%  Author:  Tharit Mohamed Hussain Khan (22058691), KCL MSc AI, 7CCSMPRJ
%
%  This file is self-contained for the listings/booktabs packages it uses.
%  To include it in the main report:  \input{week1_implementation}
%  Required preamble packages (add to main.tex if not already present):
%    \usepackage{booktabs}     % professional tables
%    \usepackage{listings}     % code listings
%    \usepackage{xcolor}       % listing colours
%    \usepackage{amsmath}      % aligned equations
% =====================================================================

% ---- Listing style (safe to keep here; remove if defined globally) ----
\lstdefinestyle{pyqvault}{
  basicstyle=\ttfamily\footnotesize,
  keywordstyle=\color{blue!70!black}\bfseries,
  commentstyle=\color{green!45!black}\itshape,
  stringstyle=\color{red!60!black},
  numberstyle=\tiny\color{gray},
  numbers=left, numbersep=6pt,
  frame=single, rulecolor=\color{gray!40},
  breaklines=true, showstringspaces=false,
  language=Python, captionpos=b,
}

% =====================================================================
\section{Foundation Layer Implementation (Week 1)}
\label{sec:impl-week1}

This section documents the foundation layer of the system: the components
built and verified in the first development sprint. The design goal of this
sprint was a fully testable pipeline skeleton --- configuration, data
contracts, dataset ingestion, retrieval, baseline policies, profiling and
scoring --- such that every later component (the retrieval gates that
constitute the project's core novelty) could be slotted into a stable,
test-covered structure rather than built against a moving target. The sprint
closed with the first real end-to-end measurement: the closed-book baseline
running on the local quantised model over 200 medical questions.

% ---------------------------------------------------------------------
\subsection{Architectural Overview}
\label{sec:impl-arch}

The system is organised as a six-layer pipeline in which each layer depends
only on the layers beneath it, communicating exclusively through the shared
data contracts of Layer~2. This strict layering is a deliberate design
decision: it allows each component to be unit-tested in isolation against
mock collaborators, and it ensures that the retrieval gates (Layer~5) can be
developed without modifying any earlier layer. The layers are:

\begin{enumerate}
  \item \textbf{Configuration} --- a four-file YAML hierarchy validated into a
        typed object;
  \item \textbf{Data contracts} --- the \texttt{UnifiedQuestion},
        \texttt{PolicyResult} and \texttt{RunRecord} dataclasses through which
        all other layers communicate;
  \item \textbf{Data ingestion} --- dataset loaders and the clinical-risk
        tagger that normalise heterogeneous sources into
        \texttt{UnifiedQuestion} objects;
  \item \textbf{Models and retrieval} --- the local LLM backend and the BM25
        lexical retriever;
  \item \textbf{Policies and gating} --- the policy abstraction and (in later
        sprints) the gate ensemble;
  \item \textbf{Evaluation} --- the profiler, the answer scorers and the
        resumable run driver.
\end{enumerate}

A consequence of this structure is that the test suite never requires the
multi-gigabyte language model: a \texttt{MockLLMBackend} that returns
deterministic text and synthetic logit arrays substitutes for the real model,
so the full suite of 61 unit tests executes in approximately half a second on
commodity hardware. This is what makes the project's verification claims
(Rubric: \emph{testing/verification}) credible and repeatable.

% ---------------------------------------------------------------------
\subsection{Configuration System}
\label{sec:impl-config}

Every experiment in this study is a point in a three-dimensional space:
hardware tier $\times$ retrieval policy $\times$ dataset. Hard-coding these
combinations would produce a combinatorial explosion of near-duplicate
scripts and would make reproducibility claims fragile. Instead, configuration
is expressed as four YAML files that are \emph{deep-merged} in increasing
order of precedence:
\[
  \texttt{base} \;\rightarrow\; \texttt{hardware} \;\rightarrow\;
  \texttt{policy} \;\rightarrow\; \texttt{experiment},
\]
so that a later file overrides only the keys it specifies, inheriting all
others. The merged dictionary is then validated into a \texttt{ProjectConfig}
Pydantic model, which rejects missing or mistyped fields \emph{before} any
expensive computation begins --- a property that matters when a single
experiment run can occupy several hours of CPU time.

A representative design decision encoded in the configuration is the
\texttt{logits\_all} flag. The entropy and margin gates require the model's
full per-token logit distribution, which the inference library only retains
when constructed with \texttt{logits\_all=True}, at a cost of roughly
256\,MB of additional memory. On the low hardware tier (4\,GB), this flag is
overridden to \texttt{false}, which automatically renders those two gates
unavailable and forces the gate ensemble to fall back to the
hallucination-probe gate alone. Encoding this trade-off declaratively in the
hardware YAML --- rather than as a runtime conditional --- means the
hardware-aware behaviour is auditable from configuration alone, supporting
the \emph{Hardware-Aware Policy Escalation} contribution.

To reconcile the terse shorthand used in experiment files (e.g.\ a bare
\texttt{policy: p1\_always\_retrieve} string, or a flat
\texttt{max\_questions} key) with the nested schema, a normalisation step
remaps these convenience keys onto their canonical nested locations prior to
validation. The configuration layer is covered by twelve unit tests
exercising the merge semantics, every hardware and policy override, and the
rejection of malformed input.

% ---------------------------------------------------------------------
\subsection{Data Contracts}
\label{sec:impl-schema}

All inter-layer communication passes through three dataclasses, isolated in a
schema module with no project-internal imports so that it can be imported
anywhere without risk of circular dependencies:

\begin{itemize}
  \item \texttt{UnifiedQuestion} --- a single medical question normalised from
        any source dataset. Multiple-choice questions carry a
        \texttt{choices} mapping and a letter \texttt{correct\_answer};
        open-ended questions carry \texttt{choices = None} and a gold answer
        string. Each question also carries a \texttt{risk\_level} assigned by
        the risk tagger.
  \item \texttt{PolicyResult} --- the output of running one policy on one
        question: the answer text, the gate decision and signal value (where
        applicable), the retrieved chunks, and profiling fields populated by
        the profiler wrapper.
  \item \texttt{RunRecord} --- a \texttt{PolicyResult} augmented with
        ground-truth scoring and experiment metadata, serialised to JSON Lines
        (JSONL) as one self-describing record per query.
\end{itemize}

The choice of an append-only JSONL audit log is central to the project's
reproducibility and resumability story. Because each line is an independent
record, an interrupted run can be resumed by reading back the set of
\texttt{question\_id}s already present in the output file and skipping them ---
essential when full experiments run for hours and may be interrupted. The
\texttt{RunRecord} additionally carries a free-form \texttt{qvault} dictionary
into which gate-specific payloads (per-token entropy arrays, the two probe
drafts, per-member votes) are written without requiring new typed columns,
keeping the schema stable as new gates are added.

% ---------------------------------------------------------------------
\subsection{Data Ingestion and Clinical-Risk Tagging}
\label{sec:impl-data}

Two dataset loaders were implemented for the primary benchmarks. The MIRAGE
loader parses the MedRAG benchmark file, which is keyed by subset
(\texttt{mmlu}, \texttt{medqa}, \texttt{medmcqa}, \texttt{pubmedqa},
\texttt{bioasq}), and emits multiple-choice \texttt{UnifiedQuestion} objects;
it tolerates both the nested subset-keyed format and a flat pre-extracted
single subset. The RAGCare-QA loader handles open-ended questions and resolves
field names flexibly (e.g.\ \texttt{question}/\texttt{query}/\texttt{input}
for the prompt), so that a HuggingFace export loads without manual remapping.
Isolating each dataset's idiosyncrasies inside its loader means that every
downstream component operates on \texttt{UnifiedQuestion} alone and is
agnostic to which dataset is running.

The clinical-risk tagger assigns each question a level in
$\{\texttt{low},\texttt{medium},\texttt{high}\}$ using a transparent,
training-free keyword heuristic with the precedence
$\textsc{high} \succ \textsc{low} \succ \textsc{medium}$ and a conservative
default of \textsc{medium}. \textsc{high}-risk cues include dosing, overdose,
drug interactions, contraindications and acute emergencies --- categories in
which an incorrect answer could directly harm a patient. The decision to use
an auditable rule set rather than a learned classifier is deliberate: in a
clinical-safety context the basis for each risk assignment must be inspectable
and defensible in the written report, and a learned classifier would both
obscure that basis and require labelled training data that does not exist for
this purpose. The tagger's limitations (notably, recall gaps on
implicitly high-risk questions) are acknowledged as a documented finding
rather than hidden. The \texttt{risk\_level} produced here is the axis along
which the \emph{Safety Envelope} analysis is later computed.

% ---------------------------------------------------------------------
\subsection{Local Model Backend and Lexical Retrieval}
\label{sec:impl-model-retrieval}

The language model is served entirely on CPU through a \texttt{LlamaBackend}
wrapping a 4-bit quantised Llama-3.2-3B model in GGUF format. Beyond ordinary
generation, the backend exposes two specialised methods that exist solely to
feed the retrieval gates:

\begin{itemize}
  \item \texttt{draft()} returns both the generated text \emph{and} the raw
        logit array of shape $[\,n_{\text{gen}} \times |V|\,]$, the direct
        input to the entropy gate;
  \item \texttt{get\_top2\_logprobs()} returns the top-two token
        log-probabilities per generated position, the input to the margin
        gate. Critically, this method relies on the completion API's
        \texttt{logprobs} output rather than the raw logit buffer, so it
        continues to function when \texttt{logits\_all=false} on the low tier.
\end{itemize}

During this sprint, \texttt{get\_top2\_logprobs()} was promoted to an abstract
method on the \texttt{LLMBackend} interface. This is a small but consequential
design choice: by placing the method on the abstraction, the forthcoming
margin gate can be written against the interface type alone, and the
\texttt{MockLLMBackend} can supply synthetic top-two log-probabilities, so the
gate is fully unit-testable without the real model. Determinism is guaranteed
by constructing the model with \texttt{temperature\,$=$\,0} (greedy decoding)
and a fixed seed; the configuration validator emits a warning if a non-zero
temperature is ever supplied, since this would silently break reproducibility.

Lexical retrieval is provided by a BM25 retriever over a corpus of
\texttt{Chunk} objects. BM25 was selected as the first retriever for two
reasons. First, it is the only retriever viable on the low hardware tier,
since it requires no embedding model, no vector index and no GPU. Second,
medical terminology is highly specific, so lexical overlap is a strong signal
over corpora such as StatPearls and medical textbooks. The retriever provides
an in-memory constructor for tests and small corpora, and a pickle-loading
constructor for the full pre-built index, and returns the top-$k$ chunks with
their BM25 scores attached. It is the lexical leg of the hybrid retriever to
be built in the next sprint, and the sole retriever used by the gated policy's
first prototype.

% ---------------------------------------------------------------------
\subsection{Baseline Policies}
\label{sec:impl-policies}

All policies implement a single abstract method,
\texttt{answer(question) -> PolicyResult}, and are constructed with the three
collaborators they may use: an LLM backend (always), a retriever (for
retrieving policies) and a gate (for the gated policy only). This uniformity
is what makes the comparative experiment tractable: the run driver invokes
\texttt{policy.answer(q)} identically for every policy, with no knowledge of
which policy it holds.

Two baseline policies were implemented and define the two ends of the
accuracy/cost spectrum against which the gated policy is later judged:

\begin{itemize}
  \item \textbf{P3 (Closed-Book)} answers purely from the model's parametric
        knowledge with no retrieval. It is the speed and energy \emph{floor}
        --- the cheapest possible policy --- and quantifies how far
        parametric-only accuracy falls short of retrieval-augmented accuracy.
  \item \textbf{P1 (Always-Retrieve)} retrieves evidence for every query and
        answers with it in context. It is the accuracy \emph{ceiling}: the
        gated policy aims to recover most of P1's accuracy while retrieving on
        a fraction of queries.
\end{itemize}

% ---------------------------------------------------------------------
\subsection{Profiling and Scoring}
\label{sec:impl-eval}

The profiler wraps any callable and returns the result alongside latency
(via a monotonic high-resolution nanosecond counter), peak memory (via
\texttt{tracemalloc}) and, optionally, energy. Energy measurement through
\texttt{codecarbon} is disabled by default and imported lazily, so that the
test suite never incurs that dependency; on real runs it is enabled and
reported in kWh. Measurement is strictly sequential: concurrent queries would
corrupt the process-level energy readings, so no parallelism is permitted in
the run loop. This is a methodological constraint imposed by the energy
metric, not an implementation limitation.

Two scorers match the two question types. The multiple-choice scorer extracts
the predicted option letter from free-form model text using a cascade of
regular expressions that handle leading letters, ``the answer is X'' phrasings
and bare capitals, then compares to the gold letter. The open-ended scorer
computes a SQuAD-style token-level $F_1$ between prediction and gold. The
token-$F_1$ function is deliberately located in the scoring module rather than
the metrics module, because the hallucination-probe gate reuses it to measure
agreement between its two drafts --- an early example of the design
anticipating the needs of the core-novelty layer.

% ---------------------------------------------------------------------
\subsection{First End-to-End Result: Closed-Book Baseline}
\label{sec:impl-result}

The sprint concluded by running the closed-book policy (P3) on the local
Llama-3.2-3B model over 200 MIRAGE questions, with full per-query profiling
and JSONL logging. This is the project's first real measurement and serves as
the parametric-knowledge floor for all subsequent comparisons.
Table~\ref{tab:p3-baseline} summarises the aggregate result, and
Table~\ref{tab:p3-byrisk} breaks accuracy down by the clinical-risk level
assigned by the tagger.

\begin{table}[ht]
  \centering
  \caption{Closed-book baseline (P3) on 200 MIRAGE questions, Llama-3.2-3B
           Q4\_K\_M, medium tier, greedy decoding, seed~42.}
  \label{tab:p3-baseline}
  \begin{tabular}{lr}
    \toprule
    \textbf{Metric} & \textbf{Value} \\
    \midrule
    Questions evaluated            & 200 \\
    Accuracy (exact letter match)  & 62.0\% \\
    Retrieval rate                 & 0\% (closed-book) \\
    Latency, median (p50)          & 7.42\,s \\
    Latency, 95th percentile (p95) & 9.25\,s \\
    Mean energy per query          & $6.77\times10^{-4}$\,kWh \\
    \bottomrule
  \end{tabular}
\end{table}

\begin{table}[ht]
  \centering
  \caption{P3 closed-book accuracy by clinical-risk level (tagger-assigned).
           Note the small $n$ in the high-risk stratum: this 200-question
           pilot is dominated by medium-risk items, and a balanced
           risk-stratified evaluation is scheduled for the full experiments.}
  \label{tab:p3-byrisk}
  \begin{tabular}{lrr}
    \toprule
    \textbf{Risk level} & \textbf{$n$} & \textbf{Accuracy} \\
    \midrule
    Low    & 28  & 64.3\% \\
    Medium & 166 & 60.8\% \\
    High   & 6   & 83.3\% \\
    \midrule
    All    & 200 & 62.0\% \\
    \bottomrule
  \end{tabular}
\end{table}

Two observations frame the later evaluation. First, a 62.0\% closed-book
accuracy establishes that the quantised 3B model retains substantial medical
knowledge parametrically, which makes the central research question --- when
retrieval can be \emph{safely skipped} --- empirically meaningful rather than
rhetorical: if closed-book accuracy were near chance, gating would be
pointless. Second, the per-query latency of roughly 7--9 seconds on CPU
quantifies the cost the gating mechanism must justify: every avoided retrieval
saves both the retrieval cost and the larger-context generation cost, and the
energy figure provides the per-query baseline against which the gate's
efficiency claims will be measured. The skew of the risk distribution in this
pilot (only six high-risk items) is itself a finding that motivates the
risk-stratified sampling planned for the full experiments.

% ---------------------------------------------------------------------
\subsection{Verification Status}
\label{sec:impl-verification}

At the close of the sprint, the foundation layer is covered by 61 unit tests,
all passing, executing in approximately 0.5\,s without any model or dataset
download by virtue of the mock backend. Coverage spans: configuration
deep-merge and override semantics (12 tests); BM25 ranking and index
round-trip (5); dataset loading and risk tagging (12); baseline policy
behaviour (3); profiler latency/memory/energy gating (6); answer scoring
across letter formats and token-$F_1$ (10); and the resumable run driver's
write-then-skip logic (2). This test-first foundation is the substrate into
which the retrieval gates --- the project's core novel contribution --- are
inserted in the following sprint.

%  Required preamble packages (add to main.tex if not already present):
%    \usepackage{booktabs}     % professional tables
%    \usepackage{listings}     % code listings
%    \usepackage{xcolor}       % listing colours
%    \usepackage{amsmath}      % aligned equations
% =====================================================================

% ---- Listing style (safe to keep here; remove if defined globally) ----
\lstdefinestyle{pyqvault}{
  basicstyle=\ttfamily\footnotesize,
  keywordstyle=\color{blue!70!black}\bfseries,
  commentstyle=\color{green!45!black}\itshape,
  stringstyle=\color{red!60!black},
  numberstyle=\tiny\color{gray},
  numbers=left, numbersep=6pt,
  frame=single, rulecolor=\color{gray!40},
  breaklines=true, showstringspaces=false,
  language=Python, captionpos=b,
}

% =====================================================================
\section{Foundation Layer Implementation (Week 1)}
\label{sec:impl-week1}

This section documents the foundation layer of the system: the components
built and verified in the first development sprint. The design goal of this
sprint was a fully testable pipeline skeleton --- configuration, data
contracts, dataset ingestion, retrieval, baseline policies, profiling and
scoring --- such that every later component (the retrieval gates that
constitute the project's core novelty) could be slotted into a stable,
test-covered structure rather than built against a moving target. The sprint
closed with the first real end-to-end measurement: the closed-book baseline
running on the local quantised model over 200 medical questions.

% ---------------------------------------------------------------------
\subsection{Architectural Overview}
\label{sec:impl-arch}

The system is organised as a six-layer pipeline in which each layer depends
only on the layers beneath it, communicating exclusively through the shared
data contracts of Layer~2. This strict layering is a deliberate design
decision: it allows each component to be unit-tested in isolation against
mock collaborators, and it ensures that the retrieval gates (Layer~5) can be
developed without modifying any earlier layer. The layers are:

\begin{enumerate}
  \item \textbf{Configuration} --- a four-file YAML hierarchy validated into a
        typed object;
  \item \textbf{Data contracts} --- the \texttt{UnifiedQuestion},
        \texttt{PolicyResult} and \texttt{RunRecord} dataclasses through which
        all other layers communicate;
  \item \textbf{Data ingestion} --- dataset loaders and the clinical-risk
        tagger that normalise heterogeneous sources into
        \texttt{UnifiedQuestion} objects;
  \item \textbf{Models and retrieval} --- the local LLM backend and the BM25
        lexical retriever;
  \item \textbf{Policies and gating} --- the policy abstraction and (in later
        sprints) the gate ensemble;
  \item \textbf{Evaluation} --- the profiler, the answer scorers and the
        resumable run driver.
\end{enumerate}

A consequence of this structure is that the test suite never requires the
multi-gigabyte language model: a \texttt{MockLLMBackend} that returns
deterministic text and synthetic logit arrays substitutes for the real model,
so the full suite of 61 unit tests executes in approximately half a second on
commodity hardware. This is what makes the project's verification claims
(Rubric: \emph{testing/verification}) credible and repeatable.

% ---------------------------------------------------------------------
\subsection{Configuration System}
\label{sec:impl-config}

Every experiment in this study is a point in a three-dimensional space:
hardware tier $\times$ retrieval policy $\times$ dataset. Hard-coding these
combinations would produce a combinatorial explosion of near-duplicate
scripts and would make reproducibility claims fragile. Instead, configuration
is expressed as four YAML files that are \emph{deep-merged} in increasing
order of precedence:
\[
  \texttt{base} \;\rightarrow\; \texttt{hardware} \;\rightarrow\;
  \texttt{policy} \;\rightarrow\; \texttt{experiment},
\]
so that a later file overrides only the keys it specifies, inheriting all
others. The merged dictionary is then validated into a \texttt{ProjectConfig}
Pydantic model, which rejects missing or mistyped fields \emph{before} any
expensive computation begins --- a property that matters when a single
experiment run can occupy several hours of CPU time.

A representative design decision encoded in the configuration is the
\texttt{logits\_all} flag. The entropy and margin gates require the model's
full per-token logit distribution, which the inference library only retains
when constructed with \texttt{logits\_all=True}, at a cost of roughly
256\,MB of additional memory. On the low hardware tier (4\,GB), this flag is
overridden to \texttt{false}, which automatically renders those two gates
unavailable and forces the gate ensemble to fall back to the
hallucination-probe gate alone. Encoding this trade-off declaratively in the
hardware YAML --- rather than as a runtime conditional --- means the
hardware-aware behaviour is auditable from configuration alone, supporting
the \emph{Hardware-Aware Policy Escalation} contribution.

To reconcile the terse shorthand used in experiment files (e.g.\ a bare
\texttt{policy: p1\_always\_retrieve} string, or a flat
\texttt{max\_questions} key) with the nested schema, a normalisation step
remaps these convenience keys onto their canonical nested locations prior to
validation. The configuration layer is covered by twelve unit tests
exercising the merge semantics, every hardware and policy override, and the
rejection of malformed input.

% ---------------------------------------------------------------------
\subsection{Data Contracts}
\label{sec:impl-schema}

All inter-layer communication passes through three dataclasses, isolated in a
schema module with no project-internal imports so that it can be imported
anywhere without risk of circular dependencies:

\begin{itemize}
  \item \texttt{UnifiedQuestion} --- a single medical question normalised from
        any source dataset. Multiple-choice questions carry a
        \texttt{choices} mapping and a letter \texttt{correct\_answer};
        open-ended questions carry \texttt{choices = None} and a gold answer
        string. Each question also carries a \texttt{risk\_level} assigned by
        the risk tagger.
  \item \texttt{PolicyResult} --- the output of running one policy on one
        question: the answer text, the gate decision and signal value (where
        applicable), the retrieved chunks, and profiling fields populated by
        the profiler wrapper.
  \item \texttt{RunRecord} --- a \texttt{PolicyResult} augmented with
        ground-truth scoring and experiment metadata, serialised to JSON Lines
        (JSONL) as one self-describing record per query.
\end{itemize}

The choice of an append-only JSONL audit log is central to the project's
reproducibility and resumability story. Because each line is an independent
record, an interrupted run can be resumed by reading back the set of
\texttt{question\_id}s already present in the output file and skipping them ---
essential when full experiments run for hours and may be interrupted. The
\texttt{RunRecord} additionally carries a free-form \texttt{qvault} dictionary
into which gate-specific payloads (per-token entropy arrays, the two probe
drafts, per-member votes) are written without requiring new typed columns,
keeping the schema stable as new gates are added.

% ---------------------------------------------------------------------
\subsection{Data Ingestion and Clinical-Risk Tagging}
\label{sec:impl-data}

Two dataset loaders were implemented for the primary benchmarks. The MIRAGE
loader parses the MedRAG benchmark file, which is keyed by subset
(\texttt{mmlu}, \texttt{medqa}, \texttt{medmcqa}, \texttt{pubmedqa},
\texttt{bioasq}), and emits multiple-choice \texttt{UnifiedQuestion} objects;
it tolerates both the nested subset-keyed format and a flat pre-extracted
single subset. The RAGCare-QA loader handles open-ended questions and resolves
field names flexibly (e.g.\ \texttt{question}/\texttt{query}/\texttt{input}
for the prompt), so that a HuggingFace export loads without manual remapping.
Isolating each dataset's idiosyncrasies inside its loader means that every
downstream component operates on \texttt{UnifiedQuestion} alone and is
agnostic to which dataset is running.

The clinical-risk tagger assigns each question a level in
$\{\texttt{low},\texttt{medium},\texttt{high}\}$ using a transparent,
training-free keyword heuristic with the precedence
$\textsc{high} \succ \textsc{low} \succ \textsc{medium}$ and a conservative
default of \textsc{medium}. \textsc{high}-risk cues include dosing, overdose,
drug interactions, contraindications and acute emergencies --- categories in
which an incorrect answer could directly harm a patient. The decision to use
an auditable rule set rather than a learned classifier is deliberate: in a
clinical-safety context the basis for each risk assignment must be inspectable
and defensible in the written report, and a learned classifier would both
obscure that basis and require labelled training data that does not exist for
this purpose. The tagger's limitations (notably, recall gaps on
implicitly high-risk questions) are acknowledged as a documented finding
rather than hidden. The \texttt{risk\_level} produced here is the axis along
which the \emph{Safety Envelope} analysis is later computed.

% ---------------------------------------------------------------------
\subsection{Local Model Backend and Lexical Retrieval}
\label{sec:impl-model-retrieval}

The language model is served entirely on CPU through a \texttt{LlamaBackend}
wrapping a 4-bit quantised Llama-3.2-3B model in GGUF format. Beyond ordinary
generation, the backend exposes two specialised methods that exist solely to
feed the retrieval gates:

\begin{itemize}
  \item \texttt{draft()} returns both the generated text \emph{and} the raw
        logit array of shape $[\,n_{\text{gen}} \times |V|\,]$, the direct
        input to the entropy gate;
  \item \texttt{get\_top2\_logprobs()} returns the top-two token
        log-probabilities per generated position, the input to the margin
        gate. Critically, this method relies on the completion API's
        \texttt{logprobs} output rather than the raw logit buffer, so it
        continues to function when \texttt{logits\_all=false} on the low tier.
\end{itemize}

During this sprint, \texttt{get\_top2\_logprobs()} was promoted to an abstract
method on the \texttt{LLMBackend} interface. This is a small but consequential
design choice: by placing the method on the abstraction, the forthcoming
margin gate can be written against the interface type alone, and the
\texttt{MockLLMBackend} can supply synthetic top-two log-probabilities, so the
gate is fully unit-testable without the real model. Determinism is guaranteed
by constructing the model with \texttt{temperature\,$=$\,0} (greedy decoding)
and a fixed seed; the configuration validator emits a warning if a non-zero
temperature is ever supplied, since this would silently break reproducibility.

Lexical retrieval is provided by a BM25 retriever over a corpus of
\texttt{Chunk} objects. BM25 was selected as the first retriever for two
reasons. First, it is the only retriever viable on the low hardware tier,
since it requires no embedding model, no vector index and no GPU. Second,
medical terminology is highly specific, so lexical overlap is a strong signal
over corpora such as StatPearls and medical textbooks. The retriever provides
an in-memory constructor for tests and small corpora, and a pickle-loading
constructor for the full pre-built index, and returns the top-$k$ chunks with
their BM25 scores attached. It is the lexical leg of the hybrid retriever to
be built in the next sprint, and the sole retriever used by the gated policy's
first prototype.

% ---------------------------------------------------------------------
\subsection{Baseline Policies}
\label{sec:impl-policies}

All policies implement a single abstract method,
\texttt{answer(question) -> PolicyResult}, and are constructed with the three
collaborators they may use: an LLM backend (always), a retriever (for
retrieving policies) and a gate (for the gated policy only). This uniformity
is what makes the comparative experiment tractable: the run driver invokes
\texttt{policy.answer(q)} identically for every policy, with no knowledge of
which policy it holds.

Two baseline policies were implemented and define the two ends of the
accuracy/cost spectrum against which the gated policy is later judged:

\begin{itemize}
  \item \textbf{P3 (Closed-Book)} answers purely from the model's parametric
        knowledge with no retrieval. It is the speed and energy \emph{floor}
        --- the cheapest possible policy --- and quantifies how far
        parametric-only accuracy falls short of retrieval-augmented accuracy.
  \item \textbf{P1 (Always-Retrieve)} retrieves evidence for every query and
        answers with it in context. It is the accuracy \emph{ceiling}: the
        gated policy aims to recover most of P1's accuracy while retrieving on
        a fraction of queries.
\end{itemize}

% ---------------------------------------------------------------------
\subsection{Profiling and Scoring}
\label{sec:impl-eval}

The profiler wraps any callable and returns the result alongside latency
(via a monotonic high-resolution nanosecond counter), peak memory (via
\texttt{tracemalloc}) and, optionally, energy. Energy measurement through
\texttt{codecarbon} is disabled by default and imported lazily, so that the
test suite never incurs that dependency; on real runs it is enabled and
reported in kWh. Measurement is strictly sequential: concurrent queries would
corrupt the process-level energy readings, so no parallelism is permitted in
the run loop. This is a methodological constraint imposed by the energy
metric, not an implementation limitation.

Two scorers match the two question types. The multiple-choice scorer extracts
the predicted option letter from free-form model text using a cascade of
regular expressions that handle leading letters, ``the answer is X'' phrasings
and bare capitals, then compares to the gold letter. The open-ended scorer
computes a SQuAD-style token-level $F_1$ between prediction and gold. The
token-$F_1$ function is deliberately located in the scoring module rather than
the metrics module, because the hallucination-probe gate reuses it to measure
agreement between its two drafts --- an early example of the design
anticipating the needs of the core-novelty layer.

% ---------------------------------------------------------------------
\subsection{First End-to-End Result: Closed-Book Baseline}
\label{sec:impl-result}

The sprint concluded by running the closed-book policy (P3) on the local
Llama-3.2-3B model over 200 MIRAGE questions, with full per-query profiling
and JSONL logging. This is the project's first real measurement and serves as
the parametric-knowledge floor for all subsequent comparisons.
Table~\ref{tab:p3-baseline} summarises the aggregate result, and
Table~\ref{tab:p3-byrisk} breaks accuracy down by the clinical-risk level
assigned by the tagger.

\begin{table}[ht]
  \centering
  \caption{Closed-book baseline (P3) on 200 MIRAGE questions, Llama-3.2-3B
           Q4\_K\_M, medium tier, greedy decoding, seed~42.}
  \label{tab:p3-baseline}
  \begin{tabular}{lr}
    \toprule
    \textbf{Metric} & \textbf{Value} \\
    \midrule
    Questions evaluated            & 200 \\
    Accuracy (exact letter match)  & 62.0\% \\
    Retrieval rate                 & 0\% (closed-book) \\
    Latency, median (p50)          & 7.42\,s \\
    Latency, 95th percentile (p95) & 9.25\,s \\
    Mean energy per query          & $6.77\times10^{-4}$\,kWh \\
    \bottomrule
  \end{tabular}
\end{table}

\begin{table}[ht]
  \centering
  \caption{P3 closed-book accuracy by clinical-risk level (tagger-assigned).
           Note the small $n$ in the high-risk stratum: this 200-question
           pilot is dominated by medium-risk items, and a balanced
           risk-stratified evaluation is scheduled for the full experiments.}
  \label{tab:p3-byrisk}
  \begin{tabular}{lrr}
    \toprule
    \textbf{Risk level} & \textbf{$n$} & \textbf{Accuracy} \\
    \midrule
    Low    & 28  & 64.3\% \\
    Medium & 166 & 60.8\% \\
    High   & 6   & 83.3\% \\
    \midrule
    All    & 200 & 62.0\% \\
    \bottomrule
  \end{tabular}
\end{table}

Two observations frame the later evaluation. First, a 62.0\% closed-book
accuracy establishes that the quantised 3B model retains substantial medical
knowledge parametrically, which makes the central research question --- when
retrieval can be \emph{safely skipped} --- empirically meaningful rather than
rhetorical: if closed-book accuracy were near chance, gating would be
pointless. Second, the per-query latency of roughly 7--9 seconds on CPU
quantifies the cost the gating mechanism must justify: every avoided retrieval
saves both the retrieval cost and the larger-context generation cost, and the
energy figure provides the per-query baseline against which the gate's
efficiency claims will be measured. The skew of the risk distribution in this
pilot (only six high-risk items) is itself a finding that motivates the
risk-stratified sampling planned for the full experiments.

% ---------------------------------------------------------------------
\subsection{Verification Status}
\label{sec:impl-verification}

At the close of the sprint, the foundation layer is covered by 61 unit tests,
all passing, executing in approximately 0.5\,s without any model or dataset
download by virtue of the mock backend. Coverage spans: configuration
deep-merge and override semantics (12 tests); BM25 ranking and index
round-trip (5); dataset loading and risk tagging (12); baseline policy
behaviour (3); profiler latency/memory/energy gating (6); answer scoring
across letter formats and token-$F_1$ (10); and the resumable run driver's
write-then-skip logic (2). This test-first foundation is the substrate into
which the retrieval gates --- the project's core novel contribution --- are
inserted in the following sprint.

%  Required preamble packages (add to main.tex if not already present):
%    \usepackage{booktabs}     % professional tables
%    \usepackage{listings}     % code listings
%    \usepackage{xcolor}       % listing colours
%    \usepackage{amsmath}      % aligned equations
% =====================================================================

% ---- Listing style (safe to keep here; remove if defined globally) ----
\lstdefinestyle{pyqvault}{
  basicstyle=\ttfamily\footnotesize,
  keywordstyle=\color{blue!70!black}\bfseries,
  commentstyle=\color{green!45!black}\itshape,
  stringstyle=\color{red!60!black},
  numberstyle=\tiny\color{gray},
  numbers=left, numbersep=6pt,
  frame=single, rulecolor=\color{gray!40},
  breaklines=true, showstringspaces=false,
  language=Python, captionpos=b,
}

% =====================================================================
\section{Foundation Layer Implementation (Week 1)}
\label{sec:impl-week1}

This section documents the foundation layer of the system: the components
built and verified in the first development sprint. The design goal of this
sprint was a fully testable pipeline skeleton --- configuration, data
contracts, dataset ingestion, retrieval, baseline policies, profiling and
scoring --- such that every later component (the retrieval gates that
constitute the project's core novelty) could be slotted into a stable,
test-covered structure rather than built against a moving target. The sprint
closed with the first real end-to-end measurement: the closed-book baseline
running on the local quantised model over 200 medical questions.

% ---------------------------------------------------------------------
\subsection{Architectural Overview}
\label{sec:impl-arch}

The system is organised as a six-layer pipeline in which each layer depends
only on the layers beneath it, communicating exclusively through the shared
data contracts of Layer~2. This strict layering is a deliberate design
decision: it allows each component to be unit-tested in isolation against
mock collaborators, and it ensures that the retrieval gates (Layer~5) can be
developed without modifying any earlier layer. The layers are:

\begin{enumerate}
  \item \textbf{Configuration} --- a four-file YAML hierarchy validated into a
        typed object;
  \item \textbf{Data contracts} --- the \texttt{UnifiedQuestion},
        \texttt{PolicyResult} and \texttt{RunRecord} dataclasses through which
        all other layers communicate;
  \item \textbf{Data ingestion} --- dataset loaders and the clinical-risk
        tagger that normalise heterogeneous sources into
        \texttt{UnifiedQuestion} objects;
  \item \textbf{Models and retrieval} --- the local LLM backend and the BM25
        lexical retriever;
  \item \textbf{Policies and gating} --- the policy abstraction and (in later
        sprints) the gate ensemble;
  \item \textbf{Evaluation} --- the profiler, the answer scorers and the
        resumable run driver.
\end{enumerate}

A consequence of this structure is that the test suite never requires the
multi-gigabyte language model: a \texttt{MockLLMBackend} that returns
deterministic text and synthetic logit arrays substitutes for the real model,
so the full suite of 61 unit tests executes in approximately half a second on
commodity hardware. This is what makes the project's verification claims
(Rubric: \emph{testing/verification}) credible and repeatable.

% ---------------------------------------------------------------------
\subsection{Configuration System}
\label{sec:impl-config}

Every experiment in this study is a point in a three-dimensional space:
hardware tier $\times$ retrieval policy $\times$ dataset. Hard-coding these
combinations would produce a combinatorial explosion of near-duplicate
scripts and would make reproducibility claims fragile. Instead, configuration
is expressed as four YAML files that are \emph{deep-merged} in increasing
order of precedence:
\[
  \texttt{base} \;\rightarrow\; \texttt{hardware} \;\rightarrow\;
  \texttt{policy} \;\rightarrow\; \texttt{experiment},
\]
so that a later file overrides only the keys it specifies, inheriting all
others. The merged dictionary is then validated into a \texttt{ProjectConfig}
Pydantic model, which rejects missing or mistyped fields \emph{before} any
expensive computation begins --- a property that matters when a single
experiment run can occupy several hours of CPU time.

A representative design decision encoded in the configuration is the
\texttt{logits\_all} flag. The entropy and margin gates require the model's
full per-token logit distribution, which the inference library only retains
when constructed with \texttt{logits\_all=True}, at a cost of roughly
256\,MB of additional memory. On the low hardware tier (4\,GB), this flag is
overridden to \texttt{false}, which automatically renders those two gates
unavailable and forces the gate ensemble to fall back to the
hallucination-probe gate alone. Encoding this trade-off declaratively in the
hardware YAML --- rather than as a runtime conditional --- means the
hardware-aware behaviour is auditable from configuration alone, supporting
the \emph{Hardware-Aware Policy Escalation} contribution.

To reconcile the terse shorthand used in experiment files (e.g.\ a bare
\texttt{policy: p1\_always\_retrieve} string, or a flat
\texttt{max\_questions} key) with the nested schema, a normalisation step
remaps these convenience keys onto their canonical nested locations prior to
validation. The configuration layer is covered by twelve unit tests
exercising the merge semantics, every hardware and policy override, and the
rejection of malformed input.

% ---------------------------------------------------------------------
\subsection{Data Contracts}
\label{sec:impl-schema}

All inter-layer communication passes through three dataclasses, isolated in a
schema module with no project-internal imports so that it can be imported
anywhere without risk of circular dependencies:

\begin{itemize}
  \item \texttt{UnifiedQuestion} --- a single medical question normalised from
        any source dataset. Multiple-choice questions carry a
        \texttt{choices} mapping and a letter \texttt{correct\_answer};
        open-ended questions carry \texttt{choices = None} and a gold answer
        string. Each question also carries a \texttt{risk\_level} assigned by
        the risk tagger.
  \item \texttt{PolicyResult} --- the output of running one policy on one
        question: the answer text, the gate decision and signal value (where
        applicable), the retrieved chunks, and profiling fields populated by
        the profiler wrapper.
  \item \texttt{RunRecord} --- a \texttt{PolicyResult} augmented with
        ground-truth scoring and experiment metadata, serialised to JSON Lines
        (JSONL) as one self-describing record per query.
\end{itemize}

The choice of an append-only JSONL audit log is central to the project's
reproducibility and resumability story. Because each line is an independent
record, an interrupted run can be resumed by reading back the set of
\texttt{question\_id}s already present in the output file and skipping them ---
essential when full experiments run for hours and may be interrupted. The
\texttt{RunRecord} additionally carries a free-form \texttt{qvault} dictionary
into which gate-specific payloads (per-token entropy arrays, the two probe
drafts, per-member votes) are written without requiring new typed columns,
keeping the schema stable as new gates are added.

% ---------------------------------------------------------------------
\subsection{Data Ingestion and Clinical-Risk Tagging}
\label{sec:impl-data}

Two dataset loaders were implemented for the primary benchmarks. The MIRAGE
loader parses the MedRAG benchmark file, which is keyed by subset
(\texttt{mmlu}, \texttt{medqa}, \texttt{medmcqa}, \texttt{pubmedqa},
\texttt{bioasq}), and emits multiple-choice \texttt{UnifiedQuestion} objects;
it tolerates both the nested subset-keyed format and a flat pre-extracted
single subset. The RAGCare-QA loader handles open-ended questions and resolves
field names flexibly (e.g.\ \texttt{question}/\texttt{query}/\texttt{input}
for the prompt), so that a HuggingFace export loads without manual remapping.
Isolating each dataset's idiosyncrasies inside its loader means that every
downstream component operates on \texttt{UnifiedQuestion} alone and is
agnostic to which dataset is running.

The clinical-risk tagger assigns each question a level in
$\{\texttt{low},\texttt{medium},\texttt{high}\}$ using a transparent,
training-free keyword heuristic with the precedence
$\textsc{high} \succ \textsc{low} \succ \textsc{medium}$ and a conservative
default of \textsc{medium}. \textsc{high}-risk cues include dosing, overdose,
drug interactions, contraindications and acute emergencies --- categories in
which an incorrect answer could directly harm a patient. The decision to use
an auditable rule set rather than a learned classifier is deliberate: in a
clinical-safety context the basis for each risk assignment must be inspectable
and defensible in the written report, and a learned classifier would both
obscure that basis and require labelled training data that does not exist for
this purpose. The tagger's limitations (notably, recall gaps on
implicitly high-risk questions) are acknowledged as a documented finding
rather than hidden. The \texttt{risk\_level} produced here is the axis along
which the \emph{Safety Envelope} analysis is later computed.

% ---------------------------------------------------------------------
\subsection{Local Model Backend and Lexical Retrieval}
\label{sec:impl-model-retrieval}

The language model is served entirely on CPU through a \texttt{LlamaBackend}
wrapping a 4-bit quantised Llama-3.2-3B model in GGUF format. Beyond ordinary
generation, the backend exposes two specialised methods that exist solely to
feed the retrieval gates:

\begin{itemize}
  \item \texttt{draft()} returns both the generated text \emph{and} the raw
        logit array of shape $[\,n_{\text{gen}} \times |V|\,]$, the direct
        input to the entropy gate;
  \item \texttt{get\_top2\_logprobs()} returns the top-two token
        log-probabilities per generated position, the input to the margin
        gate. Critically, this method relies on the completion API's
        \texttt{logprobs} output rather than the raw logit buffer, so it
        continues to function when \texttt{logits\_all=false} on the low tier.
\end{itemize}

During this sprint, \texttt{get\_top2\_logprobs()} was promoted to an abstract
method on the \texttt{LLMBackend} interface. This is a small but consequential
design choice: by placing the method on the abstraction, the forthcoming
margin gate can be written against the interface type alone, and the
\texttt{MockLLMBackend} can supply synthetic top-two log-probabilities, so the
gate is fully unit-testable without the real model. Determinism is guaranteed
by constructing the model with \texttt{temperature\,$=$\,0} (greedy decoding)
and a fixed seed; the configuration validator emits a warning if a non-zero
temperature is ever supplied, since this would silently break reproducibility.

Lexical retrieval is provided by a BM25 retriever over a corpus of
\texttt{Chunk} objects. BM25 was selected as the first retriever for two
reasons. First, it is the only retriever viable on the low hardware tier,
since it requires no embedding model, no vector index and no GPU. Second,
medical terminology is highly specific, so lexical overlap is a strong signal
over corpora such as StatPearls and medical textbooks. The retriever provides
an in-memory constructor for tests and small corpora, and a pickle-loading
constructor for the full pre-built index, and returns the top-$k$ chunks with
their BM25 scores attached. It is the lexical leg of the hybrid retriever to
be built in the next sprint, and the sole retriever used by the gated policy's
first prototype.

% ---------------------------------------------------------------------
\subsection{Baseline Policies}
\label{sec:impl-policies}

All policies implement a single abstract method,
\texttt{answer(question) -> PolicyResult}, and are constructed with the three
collaborators they may use: an LLM backend (always), a retriever (for
retrieving policies) and a gate (for the gated policy only). This uniformity
is what makes the comparative experiment tractable: the run driver invokes
\texttt{policy.answer(q)} identically for every policy, with no knowledge of
which policy it holds.

Two baseline policies were implemented and define the two ends of the
accuracy/cost spectrum against which the gated policy is later judged:

\begin{itemize}
  \item \textbf{P3 (Closed-Book)} answers purely from the model's parametric
        knowledge with no retrieval. It is the speed and energy \emph{floor}
        --- the cheapest possible policy --- and quantifies how far
        parametric-only accuracy falls short of retrieval-augmented accuracy.
  \item \textbf{P1 (Always-Retrieve)} retrieves evidence for every query and
        answers with it in context. It is the accuracy \emph{ceiling}: the
        gated policy aims to recover most of P1's accuracy while retrieving on
        a fraction of queries.
\end{itemize}

% ---------------------------------------------------------------------
\subsection{Profiling and Scoring}
\label{sec:impl-eval}

The profiler wraps any callable and returns the result alongside latency
(via a monotonic high-resolution nanosecond counter), peak memory (via
\texttt{tracemalloc}) and, optionally, energy. Energy measurement through
\texttt{codecarbon} is disabled by default and imported lazily, so that the
test suite never incurs that dependency; on real runs it is enabled and
reported in kWh. Measurement is strictly sequential: concurrent queries would
corrupt the process-level energy readings, so no parallelism is permitted in
the run loop. This is a methodological constraint imposed by the energy
metric, not an implementation limitation.

Two scorers match the two question types. The multiple-choice scorer extracts
the predicted option letter from free-form model text using a cascade of
regular expressions that handle leading letters, ``the answer is X'' phrasings
and bare capitals, then compares to the gold letter. The open-ended scorer
computes a SQuAD-style token-level $F_1$ between prediction and gold. The
token-$F_1$ function is deliberately located in the scoring module rather than
the metrics module, because the hallucination-probe gate reuses it to measure
agreement between its two drafts --- an early example of the design
anticipating the needs of the core-novelty layer.

% ---------------------------------------------------------------------
\subsection{First End-to-End Result: Closed-Book Baseline}
\label{sec:impl-result}

The sprint concluded by running the closed-book policy (P3) on the local
Llama-3.2-3B model over 200 MIRAGE questions, with full per-query profiling
and JSONL logging. This is the project's first real measurement and serves as
the parametric-knowledge floor for all subsequent comparisons.
Table~\ref{tab:p3-baseline} summarises the aggregate result, and
Table~\ref{tab:p3-byrisk} breaks accuracy down by the clinical-risk level
assigned by the tagger.

\begin{table}[ht]
  \centering
  \caption{Closed-book baseline (P3) on 200 MIRAGE questions, Llama-3.2-3B
           Q4\_K\_M, medium tier, greedy decoding, seed~42.}
  \label{tab:p3-baseline}
  \begin{tabular}{lr}
    \toprule
    \textbf{Metric} & \textbf{Value} \\
    \midrule
    Questions evaluated            & 200 \\
    Accuracy (exact letter match)  & 62.0\% \\
    Retrieval rate                 & 0\% (closed-book) \\
    Latency, median (p50)          & 7.42\,s \\
    Latency, 95th percentile (p95) & 9.25\,s \\
    Mean energy per query          & $6.77\times10^{-4}$\,kWh \\
    \bottomrule
  \end{tabular}
\end{table}

\begin{table}[ht]
  \centering
  \caption{P3 closed-book accuracy by clinical-risk level (tagger-assigned).
           Note the small $n$ in the high-risk stratum: this 200-question
           pilot is dominated by medium-risk items, and a balanced
           risk-stratified evaluation is scheduled for the full experiments.}
  \label{tab:p3-byrisk}
  \begin{tabular}{lrr}
    \toprule
    \textbf{Risk level} & \textbf{$n$} & \textbf{Accuracy} \\
    \midrule
    Low    & 28  & 64.3\% \\
    Medium & 166 & 60.8\% \\
    High   & 6   & 83.3\% \\
    \midrule
    All    & 200 & 62.0\% \\
    \bottomrule
  \end{tabular}
\end{table}

Two observations frame the later evaluation. First, a 62.0\% closed-book
accuracy establishes that the quantised 3B model retains substantial medical
knowledge parametrically, which makes the central research question --- when
retrieval can be \emph{safely skipped} --- empirically meaningful rather than
rhetorical: if closed-book accuracy were near chance, gating would be
pointless. Second, the per-query latency of roughly 7--9 seconds on CPU
quantifies the cost the gating mechanism must justify: every avoided retrieval
saves both the retrieval cost and the larger-context generation cost, and the
energy figure provides the per-query baseline against which the gate's
efficiency claims will be measured. The skew of the risk distribution in this
pilot (only six high-risk items) is itself a finding that motivates the
risk-stratified sampling planned for the full experiments.

% ---------------------------------------------------------------------
\subsection{Verification Status}
\label{sec:impl-verification}

At the close of the sprint, the foundation layer is covered by 61 unit tests,
all passing, executing in approximately 0.5\,s without any model or dataset
download by virtue of the mock backend. Coverage spans: configuration
deep-merge and override semantics (12 tests); BM25 ranking and index
round-trip (5); dataset loading and risk tagging (12); baseline policy
behaviour (3); profiler latency/memory/energy gating (6); answer scoring
across letter formats and token-$F_1$ (10); and the resumable run driver's
write-then-skip logic (2). This test-first foundation is the substrate into
which the retrieval gates --- the project's core novel contribution --- are
inserted in the following sprint.

%  Required preamble packages (add to main.tex if not already present):
%    \usepackage{booktabs}     % professional tables
%    \usepackage{listings}     % code listings
%    \usepackage{xcolor}       % listing colours
%    \usepackage{amsmath}      % aligned equations
% =====================================================================

% ---- Listing style (safe to keep here; remove if defined globally) ----
\lstdefinestyle{pyqvault}{
  basicstyle=\ttfamily\footnotesize,
  keywordstyle=\color{blue!70!black}\bfseries,
  commentstyle=\color{green!45!black}\itshape,
  stringstyle=\color{red!60!black},
  numberstyle=\tiny\color{gray},
  numbers=left, numbersep=6pt,
  frame=single, rulecolor=\color{gray!40},
  breaklines=true, showstringspaces=false,
  language=Python, captionpos=b,
}

% =====================================================================
\section{Foundation Layer Implementation (Week 1)}
\label{sec:impl-week1}

This section documents the foundation layer of the system: the components
built and verified in the first development sprint. The design goal of this
sprint was a fully testable pipeline skeleton --- configuration, data
contracts, dataset ingestion, retrieval, baseline policies, profiling and
scoring --- such that every later component (the retrieval gates that
constitute the project's core novelty) could be slotted into a stable,
test-covered structure rather than built against a moving target. The sprint
closed with the first real end-to-end measurement: the closed-book baseline
running on the local quantised model over 200 medical questions.

% ---------------------------------------------------------------------
\subsection{Architectural Overview}
\label{sec:impl-arch}

The system is organised as a six-layer pipeline in which each layer depends
only on the layers beneath it, communicating exclusively through the shared
data contracts of Layer~2. This strict layering is a deliberate design
decision: it allows each component to be unit-tested in isolation against
mock collaborators, and it ensures that the retrieval gates (Layer~5) can be
developed without modifying any earlier layer. The layers are:

\begin{enumerate}
  \item \textbf{Configuration} --- a four-file YAML hierarchy validated into a
        typed object;
  \item \textbf{Data contracts} --- the \texttt{UnifiedQuestion},
        \texttt{PolicyResult} and \texttt{RunRecord} dataclasses through which
        all other layers communicate;
  \item \textbf{Data ingestion} --- dataset loaders and the clinical-risk
        tagger that normalise heterogeneous sources into
        \texttt{UnifiedQuestion} objects;
  \item \textbf{Models and retrieval} --- the local LLM backend and the BM25
        lexical retriever;
  \item \textbf{Policies and gating} --- the policy abstraction and (in later
        sprints) the gate ensemble;
  \item \textbf{Evaluation} --- the profiler, the answer scorers and the
        resumable run driver.
\end{enumerate}

A consequence of this structure is that the test suite never requires the
multi-gigabyte language model: a \texttt{MockLLMBackend} that returns
deterministic text and synthetic logit arrays substitutes for the real model,
so the full suite of 61 unit tests executes in approximately half a second on
commodity hardware. This is what makes the project's verification claims
(Rubric: \emph{testing/verification}) credible and repeatable.

% ---------------------------------------------------------------------
\subsection{Configuration System}
\label{sec:impl-config}

Every experiment in this study is a point in a three-dimensional space:
hardware tier $\times$ retrieval policy $\times$ dataset. Hard-coding these
combinations would produce a combinatorial explosion of near-duplicate
scripts and would make reproducibility claims fragile. Instead, configuration
is expressed as four YAML files that are \emph{deep-merged} in increasing
order of precedence:
\[
  \texttt{base} \;\rightarrow\; \texttt{hardware} \;\rightarrow\;
  \texttt{policy} \;\rightarrow\; \texttt{experiment},
\]
so that a later file overrides only the keys it specifies, inheriting all
others. The merged dictionary is then validated into a \texttt{ProjectConfig}
Pydantic model, which rejects missing or mistyped fields \emph{before} any
expensive computation begins --- a property that matters when a single
experiment run can occupy several hours of CPU time.

A representative design decision encoded in the configuration is the
\texttt{logits\_all} flag. The entropy and margin gates require the model's
full per-token logit distribution, which the inference library only retains
when constructed with \texttt{logits\_all=True}, at a cost of roughly
256\,MB of additional memory. On the low hardware tier (4\,GB), this flag is
overridden to \texttt{false}, which automatically renders those two gates
unavailable and forces the gate ensemble to fall back to the
hallucination-probe gate alone. Encoding this trade-off declaratively in the
hardware YAML --- rather than as a runtime conditional --- means the
hardware-aware behaviour is auditable from configuration alone, supporting
the \emph{Hardware-Aware Policy Escalation} contribution.

To reconcile the terse shorthand used in experiment files (e.g.\ a bare
\texttt{policy: p1\_always\_retrieve} string, or a flat
\texttt{max\_questions} key) with the nested schema, a normalisation step
remaps these convenience keys onto their canonical nested locations prior to
validation. The configuration layer is covered by twelve unit tests
exercising the merge semantics, every hardware and policy override, and the
rejection of malformed input.

% ---------------------------------------------------------------------
\subsection{Data Contracts}
\label{sec:impl-schema}

All inter-layer communication passes through three dataclasses, isolated in a
schema module with no project-internal imports so that it can be imported
anywhere without risk of circular dependencies:

\begin{itemize}
  \item \texttt{UnifiedQuestion} --- a single medical question normalised from
        any source dataset. Multiple-choice questions carry a
        \texttt{choices} mapping and a letter \texttt{correct\_answer};
        open-ended questions carry \texttt{choices = None} and a gold answer
        string. Each question also carries a \texttt{risk\_level} assigned by
        the risk tagger.
  \item \texttt{PolicyResult} --- the output of running one policy on one
        question: the answer text, the gate decision and signal value (where
        applicable), the retrieved chunks, and profiling fields populated by
        the profiler wrapper.
  \item \texttt{RunRecord} --- a \texttt{PolicyResult} augmented with
        ground-truth scoring and experiment metadata, serialised to JSON Lines
        (JSONL) as one self-describing record per query.
\end{itemize}

The choice of an append-only JSONL audit log is central to the project's
reproducibility and resumability story. Because each line is an independent
record, an interrupted run can be resumed by reading back the set of
\texttt{question\_id}s already present in the output file and skipping them ---
essential when full experiments run for hours and may be interrupted. The
\texttt{RunRecord} additionally carries a free-form \texttt{qvault} dictionary
into which gate-specific payloads (per-token entropy arrays, the two probe
drafts, per-member votes) are written without requiring new typed columns,
keeping the schema stable as new gates are added.

% ---------------------------------------------------------------------
\subsection{Data Ingestion and Clinical-Risk Tagging}
\label{sec:impl-data}

Two dataset loaders were implemented for the primary benchmarks. The MIRAGE
loader parses the MedRAG benchmark file, which is keyed by subset
(\texttt{mmlu}, \texttt{medqa}, \texttt{medmcqa}, \texttt{pubmedqa},
\texttt{bioasq}), and emits multiple-choice \texttt{UnifiedQuestion} objects;
it tolerates both the nested subset-keyed format and a flat pre-extracted
single subset. The RAGCare-QA loader handles open-ended questions and resolves
field names flexibly (e.g.\ \texttt{question}/\texttt{query}/\texttt{input}
for the prompt), so that a HuggingFace export loads without manual remapping.
Isolating each dataset's idiosyncrasies inside its loader means that every
downstream component operates on \texttt{UnifiedQuestion} alone and is
agnostic to which dataset is running.

The clinical-risk tagger assigns each question a level in
$\{\texttt{low},\texttt{medium},\texttt{high}\}$ using a transparent,
training-free keyword heuristic with the precedence
$\textsc{high} \succ \textsc{low} \succ \textsc{medium}$ and a conservative
default of \textsc{medium}. \textsc{high}-risk cues include dosing, overdose,
drug interactions, contraindications and acute emergencies --- categories in
which an incorrect answer could directly harm a patient. The decision to use
an auditable rule set rather than a learned classifier is deliberate: in a
clinical-safety context the basis for each risk assignment must be inspectable
and defensible in the written report, and a learned classifier would both
obscure that basis and require labelled training data that does not exist for
this purpose. The tagger's limitations (notably, recall gaps on
implicitly high-risk questions) are acknowledged as a documented finding
rather than hidden. The \texttt{risk\_level} produced here is the axis along
which the \emph{Safety Envelope} analysis is later computed.

% ---------------------------------------------------------------------
\subsection{Local Model Backend and Lexical Retrieval}
\label{sec:impl-model-retrieval}

The language model is served entirely on CPU through a \texttt{LlamaBackend}
wrapping a 4-bit quantised Llama-3.2-3B model in GGUF format. Beyond ordinary
generation, the backend exposes two specialised methods that exist solely to
feed the retrieval gates:

\begin{itemize}
  \item \texttt{draft()} returns both the generated text \emph{and} the raw
        logit array of shape $[\,n_{\text{gen}} \times |V|\,]$, the direct
        input to the entropy gate;
  \item \texttt{get\_top2\_logprobs()} returns the top-two token
        log-probabilities per generated position, the input to the margin
        gate. Critically, this method relies on the completion API's
        \texttt{logprobs} output rather than the raw logit buffer, so it
        continues to function when \texttt{logits\_all=false} on the low tier.
\end{itemize}

During this sprint, \texttt{get\_top2\_logprobs()} was promoted to an abstract
method on the \texttt{LLMBackend} interface. This is a small but consequential
design choice: by placing the method on the abstraction, the forthcoming
margin gate can be written against the interface type alone, and the
\texttt{MockLLMBackend} can supply synthetic top-two log-probabilities, so the
gate is fully unit-testable without the real model. Determinism is guaranteed
by constructing the model with \texttt{temperature\,$=$\,0} (greedy decoding)
and a fixed seed; the configuration validator emits a warning if a non-zero
temperature is ever supplied, since this would silently break reproducibility.

Lexical retrieval is provided by a BM25 retriever over a corpus of
\texttt{Chunk} objects. BM25 was selected as the first retriever for two
reasons. First, it is the only retriever viable on the low hardware tier,
since it requires no embedding model, no vector index and no GPU. Second,
medical terminology is highly specific, so lexical overlap is a strong signal
over corpora such as StatPearls and medical textbooks. The retriever provides
an in-memory constructor for tests and small corpora, and a pickle-loading
constructor for the full pre-built index, and returns the top-$k$ chunks with
their BM25 scores attached. It is the lexical leg of the hybrid retriever to
be built in the next sprint, and the sole retriever used by the gated policy's
first prototype.

% ---------------------------------------------------------------------
\subsection{Baseline Policies}
\label{sec:impl-policies}

All policies implement a single abstract method,
\texttt{answer(question) -> PolicyResult}, and are constructed with the three
collaborators they may use: an LLM backend (always), a retriever (for
retrieving policies) and a gate (for the gated policy only). This uniformity
is what makes the comparative experiment tractable: the run driver invokes
\texttt{policy.answer(q)} identically for every policy, with no knowledge of
which policy it holds.

Two baseline policies were implemented and define the two ends of the
accuracy/cost spectrum against which the gated policy is later judged:

\begin{itemize}
  \item \textbf{P3 (Closed-Book)} answers purely from the model's parametric
        knowledge with no retrieval. It is the speed and energy \emph{floor}
        --- the cheapest possible policy --- and quantifies how far
        parametric-only accuracy falls short of retrieval-augmented accuracy.
  \item \textbf{P1 (Always-Retrieve)} retrieves evidence for every query and
        answers with it in context. It is the accuracy \emph{ceiling}: the
        gated policy aims to recover most of P1's accuracy while retrieving on
        a fraction of queries.
\end{itemize}

% ---------------------------------------------------------------------
\subsection{Profiling and Scoring}
\label{sec:impl-eval}

The profiler wraps any callable and returns the result alongside latency
(via a monotonic high-resolution nanosecond counter), peak memory (via
\texttt{tracemalloc}) and, optionally, energy. Energy measurement through
\texttt{codecarbon} is disabled by default and imported lazily, so that the
test suite never incurs that dependency; on real runs it is enabled and
reported in kWh. Measurement is strictly sequential: concurrent queries would
corrupt the process-level energy readings, so no parallelism is permitted in
the run loop. This is a methodological constraint imposed by the energy
metric, not an implementation limitation.

Two scorers match the two question types. The multiple-choice scorer extracts
the predicted option letter from free-form model text using a cascade of
regular expressions that handle leading letters, ``the answer is X'' phrasings
and bare capitals, then compares to the gold letter. The open-ended scorer
computes a SQuAD-style token-level $F_1$ between prediction and gold. The
token-$F_1$ function is deliberately located in the scoring module rather than
the metrics module, because the hallucination-probe gate reuses it to measure
agreement between its two drafts --- an early example of the design
anticipating the needs of the core-novelty layer.

% ---------------------------------------------------------------------
\subsection{First End-to-End Result: Closed-Book Baseline}
\label{sec:impl-result}

The sprint concluded by running the closed-book policy (P3) on the local
Llama-3.2-3B model over 200 MIRAGE questions, with full per-query profiling
and JSONL logging. This is the project's first real measurement and serves as
the parametric-knowledge floor for all subsequent comparisons.
Table~\ref{tab:p3-baseline} summarises the aggregate result, and
Table~\ref{tab:p3-byrisk} breaks accuracy down by the clinical-risk level
assigned by the tagger.

\begin{table}[ht]
  \centering
  \caption{Closed-book baseline (P3) on 200 MIRAGE questions, Llama-3.2-3B
           Q4\_K\_M, medium tier, greedy decoding, seed~42.}
  \label{tab:p3-baseline}
  \begin{tabular}{lr}
    \toprule
    \textbf{Metric} & \textbf{Value} \\
    \midrule
    Questions evaluated            & 200 \\
    Accuracy (exact letter match)  & 62.0\% \\
    Retrieval rate                 & 0\% (closed-book) \\
    Latency, median (p50)          & 7.42\,s \\
    Latency, 95th percentile (p95) & 9.25\,s \\
    Mean energy per query          & $6.77\times10^{-4}$\,kWh \\
    \bottomrule
  \end{tabular}
\end{table}

\begin{table}[ht]
  \centering
  \caption{P3 closed-book accuracy by clinical-risk level (tagger-assigned).
           Note the small $n$ in the high-risk stratum: this 200-question
           pilot is dominated by medium-risk items, and a balanced
           risk-stratified evaluation is scheduled for the full experiments.}
  \label{tab:p3-byrisk}
  \begin{tabular}{lrr}
    \toprule
    \textbf{Risk level} & \textbf{$n$} & \textbf{Accuracy} \\
    \midrule
    Low    & 28  & 64.3\% \\
    Medium & 166 & 60.8\% \\
    High   & 6   & 83.3\% \\
    \midrule
    All    & 200 & 62.0\% \\
    \bottomrule
  \end{tabular}
\end{table}

Two observations frame the later evaluation. First, a 62.0\% closed-book
accuracy establishes that the quantised 3B model retains substantial medical
knowledge parametrically, which makes the central research question --- when
retrieval can be \emph{safely skipped} --- empirically meaningful rather than
rhetorical: if closed-book accuracy were near chance, gating would be
pointless. Second, the per-query latency of roughly 7--9 seconds on CPU
quantifies the cost the gating mechanism must justify: every avoided retrieval
saves both the retrieval cost and the larger-context generation cost, and the
energy figure provides the per-query baseline against which the gate's
efficiency claims will be measured. The skew of the risk distribution in this
pilot (only six high-risk items) is itself a finding that motivates the
risk-stratified sampling planned for the full experiments.

% ---------------------------------------------------------------------
\subsection{Verification Status}
\label{sec:impl-verification}

At the close of the sprint, the foundation layer is covered by 61 unit tests,
all passing, executing in approximately 0.5\,s without any model or dataset
download by virtue of the mock backend. Coverage spans: configuration
deep-merge and override semantics (12 tests); BM25 ranking and index
round-trip (5); dataset loading and risk tagging (12); baseline policy
behaviour (3); profiler latency/memory/energy gating (6); answer scoring
across letter formats and token-$F_1$ (10); and the resumable run driver's
write-then-skip logic (2). This test-first foundation is the substrate into
which the retrieval gates --- the project's core novel contribution --- are
inserted in the following sprint.
